\documentclass[12pt,reqno]{amsart}

% --- Packages ---------------------------------------------------
\usepackage{amsmath,amssymb,amsthm,mathrsfs}
\usepackage{booktabs}
\usepackage[colorlinks=true,linkcolor=blue,citecolor=blue,urlcolor=blue]{hyperref}
\usepackage[margin=1.15in]{geometry}
\usepackage{enumitem}
\usepackage{microtype}


% --- Theorem environments ---------------------------------------
\newtheorem{theorem}{Theorem}[section]
\newtheorem{lemma}[theorem]{Lemma}
\newtheorem{proposition}[theorem]{Proposition}
\newtheorem{corollary}[theorem]{Corollary}
\newtheorem{conjecture}[theorem]{Conjecture}
\theoremstyle{definition}
\newtheorem{definition}[theorem]{Definition}
\newtheorem{example}[theorem]{Example}
\newtheorem{axiom}{Axiom}
\theoremstyle{remark}
\newtheorem{remark}[theorem]{Remark}
\newtheorem{question}[theorem]{Question}

% --- Notation ---------------------------------------------------
\newcommand{\ZZ}{\mathbb{Z}}
\newcommand{\QQ}{\mathbb{Q}}
\newcommand{\RR}{\mathbb{R}}
\newcommand{\CC}{\mathbb{C}}
\newcommand{\Qbar}{\overline{\QQ}}
\newcommand{\Gstar}{G^{\!*}}
\newcommand{\varpiup}{\varpi}
\newcommand{\Aut}{\operatorname{Aut}}
\newcommand{\End}{\operatorname{End}}
\newcommand{\Tr}{\operatorname{Tr}}
\newcommand{\Nm}{\operatorname{N}}
\newcommand{\agm}{\operatorname{agm}}
\newcommand{\abs}[1]{\lvert #1 \rvert}

% --- Title ------------------------------------------------------
\title[The Lemniscatic Bridge Constant]{%
  The Lemniscatic Bridge Constant $\Gstar$:\\[2pt]
  Seven Derivations and a Quadratic}

\author{}
\address{}

\subjclass[2020]{11G05, 11M06, 11F67, 33E05, 33B15}
\keywords{Lemniscate constant, fine structure constant, CM elliptic curves,
  Gamma function, Watson integral, master quadratic}

\date{\today}

\begin{document}

\begin{abstract}
We study the constant
$\Gstar = \gamma_{1}^{2}/(\sqrt{2}\,\gamma_{2}^{2}) \approx 2.9587$,
where $\gamma_{1}=\Gamma(\tfrac14)$ and $\gamma_{2}=\Gamma(\tfrac12)$,
which we call the \emph{lemniscatic bridge constant}.
We present seven equivalent representations of~$\Gstar$---from the lemniscate
arc length, the Gamma function, the complete elliptic integral, the
arithmetic-geometric mean, the Jacobi theta function at the self-dual nome,
Watson's lattice integral, and the $L$-function of the CM elliptic curve
$E\colon y^{2}=x^{3}-x$---demonstrating that $\Gstar$ sits at a nexus of
classical analysis, algebraic geometry, and number theory.

We then study the quadratic
\[
  x^{2} - 16\,{{\Gstar}}^{2}\,x + 16\,{{\Gstar}}^{3} = 0,
\]
whose coefficient~$16$ is the square of the order of the geometric
automorphism group of~$E$, and whose roots are
$x_{+}=137.036\ldots$ and $x_{-}=3.024\ldots$.
In normalised form the harmonic mean of the roots is
exactly~$2 = [\QQ(i):\QQ]$.  We conjecture that $x_{+}$ equals
$\alpha^{-1}$, the inverse fine structure constant, and discuss the open
problems that remain.

This paper is the first of three; the second~\cite{Paper2} catalogues
fifty-two representations of~$\Gstar$, and the third~\cite{Paper3} establishes
$\Gstar$ as the limit of a Wallis-type product over Fermat residue classes.
\end{abstract}

\maketitle
\tableofcontents

% ================================================================
\section{Introduction}\label{sec:intro}
% ================================================================

The fine structure constant
\begin{equation}\label{eq:alpha}
  \alpha \;=\; \frac{e^{2}}{4\pi\varepsilon_{0}\hbar c}
  \;\approx\; \frac{1}{137.036}
\end{equation}
governs the strength of the electromagnetic interaction.  Its value has been
measured to extraordinary precision---the CODATA~2022 recommended value is
$\alpha^{-1} = 137.035\,999\,177(21)$---yet no accepted theory predicts it from
first principles.

Over the decades, many authors have proposed closed-form expressions for
$\alpha^{-1}$ in terms of mathematical constants.  Wyler~\cite{Wyler} used
symmetric-space volumes to obtain $\alpha^{-1}\approx 137.036\,082$;
others~\cite{Gilmore} employed group-theoretic constructions.  These attempts
share a common weakness: the expressions are chosen \emph{ad hoc}, with no
structural reason linking the mathematical objects to the electromagnetic
coupling.

In this paper we take a different approach.  We begin with a single elliptic
curve---the unique rational curve with complex multiplication by the Gaussian
integers and $j$-invariant~$1728$---and study the quadratic polynomial whose
coefficients are intrinsic invariants of that curve.  The larger root of this
quadratic agrees with $\alpha^{-1}$ to~$1.26$~ppm.  We make no claim that this
agreement constitutes a derivation; we present it as a conjecture and identify
the precise mathematical gaps that remain.
Unlike previous attempts, the present construction is determined by the
intrinsic arithmetic of a single CM curve---the only rational elliptic curve
with $j$-invariant~$1728$ and complex multiplication by the Gaussian integers.

We will use a deliberately conservative terminology.  The constant $\Gstar$
and the master quadratic are mathematical invariants of the construction.  The
identification of the larger root with a measured electromagnetic coupling is a
separate physical readout problem.  In particular, a numerical agreement with
$\alpha^{-1}$ is not by itself a physical mechanism; a successful theory must
specify an operational map from the arithmetic/lattice structure to the
low-energy coupling.

The central object is the \emph{lemniscatic bridge constant}
\begin{equation}\label{eq:Gstar-def}
  \Gstar \;=\; \frac{\gamma_{1}^{2}}{\sqrt{2}\;\gamma_{2}^{2}}
  \;=\; 2.9586751191886388\ldots\;,
\end{equation}
where $\gamma_{1}=\Gamma(\tfrac14)$ and $\gamma_{2}=\Gamma(\tfrac12)$ are
values of the Euler Gamma function at two rational arguments.  The classical
form $\Gstar = \sqrt{2}\,\Gamma(\tfrac14)^{2}/(2\pi)$ is algebraically
equivalent, since $\pi = \gamma_{2}^{2}$.
A comprehensive catalogue of fifty-two representations of~$\Gstar$, spanning
hypergeometric evaluations, modular forms, and lattice sums, is given in the
companion paper~\cite{Paper2}.

This constant has been computed since Gauss~\cite{Gauss1799} under various
guises; we collect seven equivalent representations and prove they are all
equal (Theorem~\ref{thm:seven}).


% ================================================================
\section{The Curve and Its Invariants}\label{sec:prereq}
% ================================================================

We collect the properties of the curve $E\colon y^{2}=x^{3}-x$ that enter the
construction.  All statements are classical.

\begin{proposition}[CM structure]\label{prop:CM}
The elliptic curve $E\colon y^{2}=x^{3}-x$ over~$\QQ$ satisfies:
\begin{enumerate}[label=\textup{(\roman*)}]
  \item $j(E) = 1728$.
  \item $\End_{\Qbar}(E) \cong \ZZ[i]$
    \textup{(}complex multiplication by the Gaussian integers\textup{)}.
  \item The \emph{geometric} automorphism group is
    $\Aut_{\Qbar}(E) = \{\pm 1, \pm i\} \cong \ZZ/4\ZZ$,
    with $\abs{\Aut_{\Qbar}(E)} = 4$.
    Over~$\QQ$ itself, only $\Aut_{\QQ}(E)=\{\pm 1\}$ survives; the full
    group is defined over~$\QQ(i)$.
  \item $E(\QQ)_{\mathrm{tors}} =
    \{O,(0,0),(1,0),(-1,0)\} \cong \ZZ/2\ZZ \times \ZZ/2\ZZ$,
    with $\abs{E(\QQ)_{\mathrm{tors}}} = 4$.
  \item The conductor is $N=32$ \textup{(}LMFDB label \texttt{32.a2}\textup{)}
    and the root number is $\varepsilon(E) = +1$.
\end{enumerate}
\end{proposition}

\begin{proof}
See Silverman~\cite{Silverman}\,IV.4 for (i)--(iv) and Cremona's
tables~\cite{Cremona} for the conductor; the root number computation
follows Rohrlich~\cite{Rohrlich}.
\end{proof}

\begin{remark}\label{rem:aut-basefield}
The distinction between $\Aut_{\Qbar}(E)$ and $\Aut_{\QQ}(E)$ will matter:
the number~$16 = {|\Aut_{\Qbar}(E)|}^{2}$ that enters the master quadratic
is the square of the \emph{geometric} automorphism count.  This is the same
quantity that appears in the denominator of the BSD period ratio
(cf.~\S\ref{ssec:bsd}).
\end{remark}

The CM field is $K = \QQ(i)$, with $[K:\QQ] = 2$.

\subsection{The BSD computation}\label{ssec:bsd}

The Birch--Swinnerton-Dyer formula, proven for rank-zero CM curves by
Rubin~\cite{Rubin}, gives
\begin{equation}\label{eq:bsd}
  L(E,1) \;=\; \frac{\Omega_{+}\cdot\abs{\text{III}}\cdot
  \prod_{p}c_{p}}{%
  \abs{E(\QQ)_{\mathrm{tors}}}^{2}}
\end{equation}
where $\Omega_{+}=\varpiup$ is the real period,
$\abs{\text{III}}=1$ (Rubin~\cite{Rubin}),
$c_{2}=4$ is the Tamagawa number at $p=2$ (the unique prime of bad reduction),
$c_{p}=1$ for all odd~$p$, and $\abs{E(\QQ)_{\mathrm{tors}}}=4$.
Hence
\begin{equation}\label{eq:LE1}
  L(E,1) \;=\; \frac{\varpiup\cdot 1\cdot 4}{16}
  \;=\; \frac{\varpiup}{4} \;\neq\; 0.
\end{equation}
The nonvanishing was first proved by Coates--Wiles~\cite{CoatesWiles}.
The representation $\Gstar = 8\,L(E,1)/\gamma_{2}$
(Theorem~\ref{thm:seven},~\eqref{eq:D7}) thus links $\Gstar$ directly to
the arithmetic of~$E$ via the BSD formula.

\begin{remark}\label{rem:16inBSD}
The denominator $\abs{E(\QQ)_{\mathrm{tors}}}^{2}=16$ in~\eqref{eq:bsd}
is the same integer that will appear as the coefficient of the master quadratic.
For this specific curve, this denominator also equals
${|\Aut_{\Qbar}(E)|}^{2}$---a coincidence of small integers, not a general
identity.
\end{remark}


% ================================================================
\section{Seven Representations of $\Gstar$}\label{sec:derivations}
% ================================================================

\begin{theorem}[Equivalence of representations]\label{thm:seven}
The following seven expressions all equal~$\Gstar$:
\begin{align}
  \Gstar &= \frac{2\varpiup}{\gamma_{2}}
    &&\textup{(Lemniscate)}\label{eq:D1}\\[4pt]
  \Gstar &= \frac{\gamma_{1}^{2}}{\sqrt{2}\,\gamma_{2}^{2}}
    &&\textup{(Gamma function, $\pi$-free)}\label{eq:D2}\\[4pt]
  \Gstar &= \frac{2\sqrt{2}}{\gamma_{2}}\,K\!\bigl(\tfrac{1}{\sqrt{2}}\bigr)
    &&\textup{(Elliptic integral)}\label{eq:D3}\\[4pt]
  \Gstar &= \frac{2\gamma_{2}}{\agm(1,\sqrt{2})}
    &&\textup{(Arithmetic-geometric mean)}\label{eq:D4}\\[4pt]
  \Gstar &= \gamma_{2}\sqrt{2}\;\theta_{3}(e^{-\gamma_{2}^{2}})^{2}
    &&\textup{(Theta function at self-dual nome)}\label{eq:D5}\\[4pt]
  \Gstar &= \sqrt{2\gamma_{2}^{2}\,W_{3}}
    &&\textup{(Watson integral)}\label{eq:D6}\\[4pt]
  \Gstar &= \frac{8\,L(E,1)}{\gamma_{2}}
    &&\textup{($L$-function of $E$)}\label{eq:D7}
\end{align}
Here $\varpiup = \gamma_{1}^{2}/(2\sqrt{2}\,\gamma_{2})$ is the lemniscate
constant, $K(k)$ the complete elliptic integral of the first kind,
$\theta_{3}(q) = \sum_{n\in\ZZ}q^{n^{2}}$ the Jacobi theta function, and
$W_{3} = \gamma_{1}^{4}/(4\gamma_{2}^{6})$ the Watson integral for the BCC
lattice \textup{(Watson~\cite{Watson}; see also
Zucker--Joyce~\cite{ZuckerJoyce})}.
An extended catalogue of forty-five further representations is given
in~\cite{Paper2}.
\end{theorem}

\begin{proof}
Each reduces to~\eqref{eq:D2} by a single classical identity:
\begin{itemize}
  \item \eqref{eq:D1}: $\varpiup = \Gstar\gamma_{2}/2$ by definition.
  \item \eqref{eq:D3}: Legendre~\cite{Legendre} evaluates
    $K(1/\!\sqrt{2}) = \gamma_{1}^{2}/(4\gamma_{2})$.
  \item \eqref{eq:D4}: Gauss~\cite{Gauss1799} proved
    $\agm(1,\sqrt{2}) = \gamma_{2}^{2}/\varpiup$
    (see also Borwein--Borwein~\cite{BorweinBorwein}).
  \item \eqref{eq:D5}: The Chowla--Selberg formula~\cite{ChowlaSelberg}
    gives
    $\theta_{3}(e^{-\gamma_{2}^{2}})
    = \gamma_{1}/(\gamma_{2}^{3/2}\cdot 2^{1/4})$.
  \item \eqref{eq:D6}: Watson~\cite{Watson} proved
    $W_{3} = \gamma_{1}^{4}/(4\gamma_{2}^{6})$; then
    $2\gamma_{2}^{2}W_{3} = \gamma_{1}^{4}/(2\gamma_{2}^{4}) = {{\Gstar}}^{2}$.
  \item \eqref{eq:D7}: $L(E,1)=\varpiup/4$ by~\eqref{eq:LE1}.\qedhere
\end{itemize}
\end{proof}


% ================================================================
\section{The Triad Identity}\label{sec:triad}
% ================================================================

\begin{theorem}[Triad invariant]\label{thm:triad}
$\pi = 4\,\varpiup^{2}/{{\Gstar}}^{2}$.
\end{theorem}

\begin{proof}
$\Gstar = 2\varpiup/\gamma_{2}$ and $\pi=\gamma_{2}^{2}$.
\end{proof}


% ================================================================
\section{The Master Quadratic}\label{sec:quadratic}
% ================================================================

We now state the central conjectural construction.

\begin{conjecture}[Master quadratic]\label{conj:master}
The polynomial
\begin{equation}\label{eq:master-quad}
  P(x) \;=\; x^{2} - 16\,{{\Gstar}}^{2}\,x + 16\,{{\Gstar}}^{3}
\end{equation}
is the characteristic polynomial of a canonically defined self-pairing on the
CM curve $E\colon y^{2}=x^{3}-x$, with $16={|\Aut_{\Qbar}(E)|}^{2}$
and~$\Gstar$ the scaled period.
\end{conjecture}

The polynomial~$P$ is unambiguously defined: its coefficients are built from
$\Gstar$ and the integer~$16$.  What is conjectural is the claim that~$P$
arises canonically from the geometry of~$E$.  We discuss the evidence for each
ingredient.

\subsection{Why degree~$2$?}\label{ssec:degree}

The CM field $K=\QQ(i)$ has $[K:\QQ]=2$.  Every endomorphism $\varphi\in\End(E)$
satisfies a degree-$2$ characteristic polynomial $t^{2}-\Tr(\varphi)\,t+\deg(\varphi)=0$
over~$\ZZ$.  It is therefore natural---though not logically forced---that a
``global'' coupling polynomial should also have degree~$2$.

\begin{question}\label{q:degree}
Is there a rigorous argument showing that the self-consistent coupling of~$E$
on~$\ZZ^{3}$ must satisfy a polynomial of degree exactly~$[K:\QQ]$?
\end{question}

\subsection{Why coefficient~$16$?}\label{ssec:coeff}

The integer~$16 = {|\Aut_{\Qbar}(E)|}^{2} = \abs{E(\QQ)_{\mathrm{tors}}}^{2}$
appears in the BSD formula~\eqref{eq:bsd} as the denominator of the period ratio.
It is the natural ``normalisation divisor'' for this curve.  We conjecture that
it enters the master quadratic for the same structural reason.

\begin{question}\label{q:coeff}
Does the BSD denominator $\abs{E(\QQ)_{\mathrm{tors}}}^{2}$ canonically
determine the coefficient of a self-consistency equation on~$E$?
\end{question}

\subsection{Why Sum~$=$~Product?}\label{ssec:sep}

In the normalised variable $u = x/\Gstar$, the quadratic becomes
\begin{equation}\label{eq:norm-quad}
  u^{2} - S\,u + S = 0, \qquad S = 16\,\Gstar.
\end{equation}
The distinctive feature is that the linear coefficient and the constant term
are \emph{equal}.  By Vieta:
\begin{equation}\label{eq:vieta}
  u_{+}+u_{-} \;=\; u_{+}\cdot u_{-} \;=\; S.
\end{equation}

\begin{theorem}[Harmonic mean]\label{thm:harmonic}
For any quadratic $u^{2}-Su+S=0$ with $S\neq 0$, the harmonic mean of the
roots is
\begin{equation}\label{eq:harmonic}
  H \;=\; \frac{2\,u_{+}u_{-}}{u_{+}+u_{-}}
  \;=\; \frac{2S}{S} \;=\; 2.
\end{equation}
\end{theorem}

\begin{proof}
Immediate from~\eqref{eq:vieta}.
\end{proof}

The coincidence $H = 2 = [K:\QQ]$ is the strongest structural evidence
for~\eqref{eq:master-quad}: if one accepts that the polynomial should be
degree~$2$ (matching $[K:\QQ]$) and should have Sum~$=$~Product (giving
$H=[K:\QQ]$), then~\eqref{eq:master-quad} is the unique result.

\begin{conjecture}[Self-duality forces Sum~$=$~Product]\label{conj:TrN}
The root number $\varepsilon(E)=+1$ makes the functional equation
$\Lambda(E,s)=\Lambda(E,2-s)$ an identity at $s=1$.  We conjecture that this
self-duality, when applied to a suitably defined ``global coupling operator''
on the Galois eigenspaces of~$T_{0}(E)$, forces the operator's trace and
determinant to be equal.
\end{conjecture}

\begin{remark}\label{rem:gap}
Conjecture~\ref{conj:TrN} is the central open problem.
The ``global coupling operator'' is not yet rigorously defined.  Candidate
constructions include the Petersson self-pairing of the weight-$2$ newform
$f\in S_{2}(\Gamma_{0}(32))$ associated to~$E$, decomposed by the
$\operatorname{Gal}(K/\QQ)$ action.  We do not pursue this further here.
\end{remark}


\subsection{Roots}\label{ssec:roots}

\begin{theorem}[Explicit roots]\label{thm:roots}
The roots of~\eqref{eq:master-quad} are
\begin{align}
  x_{+} &= 8\,{{\Gstar}}^{2} + 4\,\Gstar\sqrt{\Gstar(4\Gstar-1)}
  = 137.0361714582\ldots\,,\label{eq:xp}\\[3pt]
  x_{-} &= 8\,{{\Gstar}}^{2} - 4\,\Gstar\sqrt{\Gstar(4\Gstar-1)}
  = 3.0239639163\ldots\,.\label{eq:xm}
\end{align}
The discriminant is $\Delta = 64\,{{\Gstar}}^{3}(4\Gstar-1)>0$.
Both roots are transcendental: since $\gamma_{1}$ and
$\gamma_{2}^{2}=\pi$ are algebraically independent over~$\QQ$
\textup{(Nesterenko~\cite{Nesterenko})}, $\Gstar$ is transcendental, and a
quadratic in~$\Gstar$ with integer-ratio structure cannot produce algebraic
roots.
\end{theorem}

\begin{proof}
Quadratic formula.  The transcendence claim follows from
Nesterenko~\cite{Nesterenko}, Theorem~1.
\end{proof}

\begin{corollary}\label{cor:floor}
$\lfloor x_{-}\rfloor = 3$.
\end{corollary}

\subsection{The $k$-family}\label{ssec:kfamily}

Replacing~$16$ by a parameter~$k$, the normalised quadratic
$u^{2} - k\Gstar\,u + k\Gstar = 0$
still has $H=2$ for every $k>0$.  The discriminant vanishes at
$k_{\mathrm{crit}} = 4/\Gstar \approx 1.352$; below this value the roots are
complex.


% ================================================================
\section{The Dual Convergence}\label{sec:dual}
% ================================================================

\begin{theorem}[Dual reading of~$r_{2}$]\label{thm:dual}
Let $r_{2}(n) = \#\{(a,b)\in\ZZ^{2}: a^{2}+b^{2}=n\}$ and
$R(N)=\sum_{n=0}^{N}r_{2}(n)$.  Then:
\begin{enumerate}[label=\textup{(\roman*)}]
  \item $R(N)/N \to \pi$ as $N\to\infty$
    \quad \textup{(Gauss circle theorem)}.
  \item $\sum_{n=0}^{\infty}r_{2}(n)\,e^{-\gamma_{2}^{2}n}
    = \Gstar/(\gamma_{2}\sqrt{2})$
    \quad \textup{(Chowla--Selberg)}.
\end{enumerate}
The same function~$r_{2}$ produces~$\pi$ by unweighted averaging and~$\Gstar$
by exponential weighting at the self-dual nome.
\end{theorem}


% ================================================================
\section{The Gamma-Primitive Basis}\label{sec:gamma}
% ================================================================

\begin{theorem}[Two-primitive basis]\label{thm:basis}
Every object in the elliptic hierarchy at the lemniscatic modulus
$k=1/\!\sqrt{2}$ is a rational expression in~$\gamma_{1}$, $\gamma_{2}$,
and~$\sqrt{2}$:
\begin{alignat}{2}
  \varpiup &= \frac{\gamma_{1}^{2}}{2\sqrt{2}\,\gamma_{2}}\,,
    &\qquad K(1/\!\sqrt{2}) &= \frac{\gamma_{1}^{2}}{4\gamma_{2}}\,,
    \label{eq:basis2}\\[4pt]
  \agm(1,\sqrt{2}) &= \frac{2\gamma_{2}^{3}}{\gamma_{1}^{2}}\,,
    &\qquad E(1/\!\sqrt{2})
    &= \frac{8\gamma_{2}^{4}+\gamma_{1}^{4}}{8\gamma_{1}^{2}\gamma_{2}}\,.
    \label{eq:basis3}
\end{alignat}
\end{theorem}

\begin{proof}
Legendre evaluations at $k=1/\!\sqrt{2}$
(cf.\ Borwein--Borwein~\cite{BorweinBorwein}).
The $E$~form follows from the Legendre relation: at $k=1/\!\sqrt{2}$ we have
$K=K'$ and $E=E'$, so $2KE-K^{2}=\gamma_{2}^{2}/2$.
\end{proof}

\begin{corollary}\label{cor:elliptic}
At $k=1/\!\sqrt{2}$:
\textup{(i)}~$E = (K + \agm(1,\sqrt{2}))/2$;
\textup{(ii)}~$E/K = 1/2 + 2/{{\Gstar}}^{2}$;
\textup{(iii)}~$\Pi(1/2, 1/\!\sqrt{2}) = 2E$.
\end{corollary}


% ================================================================
\section{The Conjecture}\label{sec:physics}
% ================================================================

\begin{conjecture}[Lattice-CM]\label{conj:lattice}
The inverse fine structure constant equals the larger root of the master
quadratic:
\[
  \alpha^{-1} \;=\; x_{+} \;=\;
  8\,{{\Gstar}}^{2} + 4\,\Gstar\sqrt{\Gstar(4\Gstar-1)}\,.
\]
\end{conjecture}

\noindent\textbf{Numerical evidence.}
$x_{+}=137.0362\ldots$ vs.\ CODATA $\alpha^{-1}=137.0360(2)$.
Agreement: $1.26$~ppm, with zero adjustable parameters in the leading term.

\smallskip\noindent\textbf{Structural observations.}
\begin{itemize}
  \item $\lfloor x_{-}\rfloor = 3$, the number of spatial dimensions and of
    quark colour charges.
  \item $D=3$ is the unique dimension $D\in\{1,\ldots,6\}$ for which the
    analogous quadratic (with the $D$-dimensional Watson integral) satisfies
    $\lfloor x_{-}^{(D)}\rfloor = D$.
\end{itemize}

\begin{remark}[Running of~$\alpha$]\label{rem:running}
The electromagnetic coupling runs with energy scale; $\alpha^{-1}\approx 137$
is the low-energy (infrared) value.  If Conjecture~\ref{conj:lattice} is
correct, $x_{+}$ should be identified with this infrared limit.  The
connection to the renormalization group flow is an open problem.
\end{remark}

\begin{remark}[Comparison with prior work]\label{rem:wyler}
Wyler~\cite{Wyler} obtained $\alpha^{-1}\approx 137.036\,08$ from
symmetric-space volumes; the expression has no structural link to an elliptic
curve and the agreement is weaker.
The present approach differs in that all ingredients---$\Gstar$, $16$, the
degree~$2$---are intrinsic invariants of a single CM curve.  Whether this
structural coherence elevates the result beyond coincidence remains to be
established.
\end{remark}

\begin{remark}[Precision formula]\label{rem:precision}
We note that the expansion
$\alpha^{-1} = x_{+} - \tfrac{9}{47}\abs{\varepsilon}
+ \tfrac{5}{64}\abs{\varepsilon}^{2} - \cdots$
with $\varepsilon = e^{\pi}-\pi-20\approx -9\times 10^{-4}$
matches CODATA to many digits.  However, the rational coefficients in this
expansion are \emph{observed}, not derived, and the expansion variable is
chosen because $e^{\pi}$ is a famous near-integer.  We include this observation
for completeness; it should not be taken as evidence for
Conjecture~\ref{conj:lattice}.
\end{remark}


% ================================================================
\section{Open Problems}\label{sec:open}
% ================================================================

\begin{enumerate}[label=(\arabic*)]
  \item \textbf{The global coupling operator.}
    Does there exist a canonically defined operator on the motive of~$E$ whose
    characteristic polynomial is~$P(x)$ as in~\eqref{eq:master-quad}?
    (See Conjecture~\ref{conj:master}.)

  \item \textbf{Why Sum~$=$~Product?}
    Does the self-duality $\Lambda(E,s)=\Lambda(E,2-s)$ at $s=1$ force the
    trace and determinant of a coupling operator on~$E$ to be equal?
    (See Conjecture~\ref{conj:TrN}.)

  \item \textbf{Why~$16$?}
    Is there a functorial construction that selects
    $16 = {|\Aut_{\Qbar}(E)|}^{2}$ as the coefficient of both the BSD period
    ratio and the master quadratic, or is the coincidence accidental?

  \item \textbf{Physical mechanism.}
    Can lattice gauge theory on~$\ZZ^{3}$, or any other physical framework,
    produce Conjecture~\ref{conj:lattice} as a consequence of a dynamical
    principle?

  \item \textbf{Operational readout.}
    Can one define a finite, non-circular observable whose value is the
    interaction strength measured as the low-energy electromagnetic coupling,
    without inserting $\alpha$ or an equivalent calibration into the definition?

  \item \textbf{Running coupling.}
    If $x_{+}=\alpha^{-1}$ at zero momentum transfer, what determines the
    beta function $\beta(\alpha)$ within this framework?  A complete theory
    must predict the running, not just the infrared fixed point.
\end{enumerate}


% ================================================================
\section{Conclusion}\label{sec:conclusion}
% ================================================================

The lemniscatic bridge constant
$\Gstar = \gamma_{1}^{2}/(\sqrt{2}\,\gamma_{2}^{2})$
has been computed since Gauss.  It admits seven equivalent representations
spanning classical analysis, lattice combinatorics, and arithmetic geometry;
a further forty-five appear in the companion catalogue~\cite{Paper2}, and a
Wallis-type product formula is established in~\cite{Paper3}.

The quadratic $x^{2}-16{{\Gstar}}^{2}x+16{{\Gstar}}^{3}=0$, with coefficient
$16={|\Aut_{\Qbar}(E)|}^{2}$ and normalised harmonic mean
$H=2=[\QQ(i):\QQ]$, has a larger root of~$137.036\ldots$  Whether this root
is the inverse fine structure constant is a conjecture, and more specifically
an unsolved readout problem rather than a theorem of the present paper.  We
have identified three precise mathematical questions (the global coupling
operator, the Sum~$=$~Product condition, and the role of~$16$), together with
the physical readout and running-coupling problems, whose resolution would
settle the matter.

We leave the question open and offer the bridge for crossing.


% ================================================================
% REFERENCES
% ================================================================

\begin{thebibliography}{20}

\bibitem{BorweinBorwein}
J.~M.~Borwein and P.~B.~Borwein,
\emph{Pi and the AGM},
Wiley, New York, 1987.

\bibitem{ChowlaSelberg}
S.~Chowla and A.~Selberg,
\emph{On Epstein's zeta-function},
J.~Reine Angew.\ Math.\ \textbf{227} (1967), 86--110.

\bibitem{Cremona}
J.~E.~Cremona,
\emph{Algorithms for Modular Elliptic Curves}, 2nd ed.,
Cambridge Univ.\ Press, 1997.
See also \url{https://www.lmfdb.org/EllipticCurve/Q/32/a/2}.

\bibitem{CoatesWiles}
J.~Coates and A.~Wiles,
\emph{On the conjecture of Birch and Swinnerton-Dyer},
Invent.\ Math.\ \textbf{39} (1977), 223--251.

\bibitem{Gauss1799}
C.~F.~Gauss,
\emph{Arithmetisch-geometrisches Mittel} (unpublished, 1799); see
\emph{Werke}, Bd.~III, 361--403.

\bibitem{Gilmore}
R.~Gilmore,
\emph{The fine structure constant},
unpublished note, Drexel University, 2004.

\bibitem{Legendre}
A.-M.~Legendre,
\emph{Exercices de calcul int\'{e}gral}, Vol.~1, Paris, 1811.

\bibitem{Nesterenko}
Yu.~V.~Nesterenko,
\emph{Modular functions and transcendence questions},
Mat.\ Sb.\ \textbf{187} (1996), 65--96.

\bibitem{Paper2}
\emph{Fifty-Two Faces of the Lemniscatic Bridge Constant},
companion paper (forthcoming).

\bibitem{Paper3}
\emph{A Wallis-Type Product for $G^{\!*}$ over Fermat Residue Classes},
companion paper (forthcoming).

\bibitem{Rohrlich}
D.~E.~Rohrlich,
\emph{Root numbers}, in: Arithmetic of $L$-functions, IAS/Park City Math.\
Ser.\ \textbf{18}, AMS, 2011, 353--448.

\bibitem{Rubin}
K.~Rubin,
\emph{The ``main conjectures'' of Iwasawa theory for imaginary quadratic
fields},
Invent.\ Math.\ \textbf{103} (1991), 25--68.

\bibitem{Silverman}
J.~H.~Silverman,
\emph{The Arithmetic of Elliptic Curves}, 2nd ed., Grad.\ Texts in Math.\
\textbf{106}, Springer, 2009.

\bibitem{Watson}
G.~N.~Watson,
\emph{Three triple integrals},
Quart.\ J.\ Math.\ Oxford \textbf{10} (1939), 266--276.

\bibitem{Wyler}
A.~Wyler,
\emph{L'espace sym\'etrique du groupe des \'equations de Maxwell},
C.~R.\ Acad.\ Sci.\ Paris \textbf{269A} (1969), 743--745.

\bibitem{ZuckerJoyce}
I.~J.~Zucker and G.~S.~Joyce,
\emph{Special values of the hypergeometric series~II},
Math.\ Proc.\ Cambridge Philos.\ Soc.\ \textbf{131} (2001), 309--319.

\end{thebibliography}

\end{document}
